\chapter{Milne Theorem 3.2: Rational maps into abelian varieties}
\label{chap:Albanese_Thm32RationalMapExtension}
% archon:covers AlgebraicJacobian/Albanese/Thm32RationalMapExtension.lean

\section{Setup and statement}
\label{sec:thm32_setup}

Fix an algebraically closed field \(\bar k\), a nonsingular variety \(X\) over
\(\bar k\), and an abelian variety \(A\) over \(\bar k\).  Thus \(X\) is
integral and smooth, while \(A\) is a complete smooth group variety.  Following
Milne, the extension statement is the one in \cite{milne-av}.  A
rational map \(f \colon X \dashrightarrow A\) is represented by a regular
morphism \(\varphi_U \colon U \to A\) on a dense open subset
\(U \subseteq X\), two such pairs \((U, \varphi_U)\) and \((U', \varphi_{U'})\) being equivalent
if \(\varphi_U\) and \(\varphi_{U'}\) agree on \(U \cap U'\).

\begin{theorem}[Milne Theorem~3.2: a rational map from a nonsingular variety to an abelian variety extends]\leanok
  \label{thm:rational_map_to_av_extends}
  \lean{AlgebraicGeometry.Scheme.RationalMap.extend_to_av}
  \dcref{custom}

  Let \(X\) be a nonsingular variety over an algebraically closed field \(\bar k\), let \(A\) be
  an abelian variety over \(\bar k\), and let \(f \colon X \dashrightarrow A\) be a rational
  map. Then \(f\) is defined on the whole of \(X\): there is a unique regular morphism
  \(\widetilde f \colon X \to A\) whose restriction to some (equivalently, every) dense open
  set of definition \(U \subseteq X\) of \(f\) agrees with \(\varphi_U\).
\end{theorem}

\begin{proof}
  \uses{lem:av_indeterminacyLocus_eq_empty}
  Let \(U\) be the maximal open subset on which \(f\) is defined, and put
  \(Z=X\setminus U\).  By \cref{lem:av_indeterminacyLocus_eq_empty},
  \(Z=\varnothing\).  Hence \(U=X\), and the representative of \(f\) on
  \(U\) is the desired morphism \(\widetilde f\colon X\to A\).

  If two morphisms \(X\to A\) represent \(f\), they agree on a dense open
  subset of \(X\).  Their equalizer is closed because \(A\) is separated,
  and it contains that dense open subset.  Since \(X\) is reduced, the two
  morphisms are equal.  Thus the extension is unique.
\end{proof}

\begin{remark}
  \label{rmk:thm32_role_of_ab}
  \dcref{custom}
  The Auslander--Buchsbaum theorem is not needed for
  \cref{thm:rational_map_to_av_extends}.  The codimension-one argument uses
  regularity only to obtain normality and discrete valuation rings at prime
  divisors; properness then supplies extensions by the valuative criterion.
  It does not extend across a codimension-two locus by a depth argument.
  The stronger Cohen--Macaulay consequences of Auslander--Buchsbaum in
  \cref{chap:Albanese_AuslanderBuchsbaum} describe the same regular local
  rings but are logically separate from Milne's proof.
\end{remark}

\section{Why the proof is characteristic-free}
\label{sec:thm32_char_free}

Both ingredients are independent of the characteristic:
\begin{itemize}
  \item \emph{Theorem~3.1 / codim-\(\geq 2\) extension} (\cref{thm:codim_one_extension}):
    a codimension-one local ring on a normal variety is a discrete valuation
    ring, and the valuative criterion of properness extends the map over that
    local ring.
  \item \emph{Lemma~3.3 / pure-codim-\(1\) indeterminacy for maps into a group variety}
    (\cref{lem:milne_codim1_indeterminacy}): Milne's proof forms the difference map
    \(\Phi(x, y) := \varphi(x)\,\varphi(y)^{-1} \colon X \times X \dashrightarrow A\) and
    examines its pole divisors along the diagonal.  This uses the group law,
    regularity, and the codimension-one support of principal divisors, none
    of which imposes a characteristic restriction.
\end{itemize}
Consequently \cref{thm:rational_map_to_av_extends} holds over every algebraically closed
field, in arbitrary characteristic.

\section{Application to the symmetric-product Albanese map}
\label{sec:thm32_application_to_a4d}

Let \(C\) be a complete nonsingular curve of genus \(g>0\), let \(J\) be
its Jacobian, and let \(\varphi\colon C\to A\) send a chosen point \(P\) to
zero.  The symmetric map
\[
  (P_1,\ldots,P_g)\longmapsto
  \varphi(P_1)+\cdots+\varphi(P_g)
\]
descends to a morphism \(\varphi^{(g)}\colon C^{(g)}\to A\).  The
Abel--Jacobi morphism \(C^{(g)}\to J\) is birational, so composing
\(\varphi^{(g)}\) with its rational inverse gives a rational map
\(J\dashrightarrow A\).  Since \(J\) is nonsingular,
\cref{thm:rational_map_to_av_extends} makes this map regular.  It factors
\(\varphi\) through \(C\to J\), and the rigidity theorem shows that it is a
homomorphism.  This is the existence step in the Albanese universal property
of \(J\) proved in \cref{chap:Albanese_AlbaneseUP}.

\section{Structure of the proof}
\label{sec:thm32_lean_encoding}

The argument separates into three elementary reductions.  Smoothness makes
the abelian variety reduced; geometric irreducibility then makes it
integral.  Milne's two indeterminacy results force the indeterminacy locus to
be empty.  Finally, a rational map with empty indeterminacy locus has a
global representative, unique because its source is reduced and its target
is separated.

\section{Auxiliary lemmas}
\label{sec:thm32_lean_helpers}

We record the reductions concerning the target and the indeterminacy locus.

\begin{lemma}[A scheme smooth over a field is reduced]\leanok
  \label{lem:isReduced_of_smooth_over_field}
  \dcref{custom}
  \lean{AlgebraicGeometry.Scheme.RationalMap.isReduced_of_smooth_over_field}
  Let \(A\) be a scheme over an algebraically closed field \(\bar k\).  If
  the structure morphism \(A\to\Spec\bar k\) is smooth, then \(A\) is
  reduced.
\end{lemma}
\begin{proof}
  An algebraically closed field is perfect.  A scheme smooth over a
  perfect field is regular, since its local rings are geometrically
  regular over the field.  Regular local rings are reduced, and
  reducedness is local; hence \(A\) is reduced.
\end{proof}

\begin{lemma}[Integrality of the abelian variety]\leanok
  \label{lem:av_isIntegral_of_smooth_geomIrred}
  \dcref{custom}
  \lean{AlgebraicGeometry.Scheme.RationalMap.av_isIntegral_of_smooth_geomIrred}
  An abelian variety \(A\) over an algebraically closed field \(\bar k\) is
  integral.  Indeed, it is geometrically irreducible and smooth over
  \(\bar k\), hence irreducible and reduced.
\end{lemma}
\begin{proof}
  \uses{lem:isReduced_of_smooth_over_field}
  The spectrum of \(\bar k\) consists of one point.  Geometric
  irreducibility of the structure morphism therefore implies that the
  underlying topological space of \(A\) is irreducible.  By
  \cref{lem:isReduced_of_smooth_over_field}, \(A\) is also reduced.  A
  nonempty scheme is integral precisely when it is irreducible and
  reduced, so \(A\) is integral.
\end{proof}

\begin{lemma}[Emptiness of the indeterminacy locus of a rational map into an abelian variety]\leanok
  \label{lem:av_indeterminacyLocus_eq_empty}
  \dcref{custom}
  \lean{AlgebraicGeometry.Scheme.RationalMap.av_indeterminacyLocus_eq_empty}
  For a rational map \(f \colon X \dashrightarrow A\) from a nonsingular variety
  \(X\) to an abelian variety \(A\) over an algebraically closed field \(\bar k\),
  the indeterminacy locus \(Z(f)\) is empty.
\end{lemma}
\begin{proof}
  \uses{lem:av_isIntegral_of_smooth_geomIrred, lem:milne_codim1_indeterminacy, thm:codim_one_extension}
  The abelian variety \(A\) is integral
  (\cref{lem:av_isIntegral_of_smooth_geomIrred}), so Milne's Lemma~3.3
  (\cref{lem:milne_codim1_indeterminacy}) applies and gives the dichotomy: either
  \(Z(f) = \emptyset\), or every point of \(Z(f)\) specialises from a point
  \(z \in Z(f)\) of coheight \(1\). In the second branch, applying the
  codim-\(\geq 2\) indeterminacy theorem (\cref{thm:codim_one_extension},
  Milne~3.1) to such a \(z \in Z(f)\) forces \(\mathrm{coheight}(z) \geq 2\),
  contradicting \(\mathrm{coheight}(z) = 1\). Hence \(Z(f) = \emptyset\).
\end{proof}

\section{Relations with neighbouring results}
\label{sec:thm32_out_of_scope}

Milne's codimension estimate and pure-codimension-one theorem are proved in
\cref{chap:Albanese_CodimOneExtension} as
\cref{thm:codim_one_extension} and
\cref{lem:milne_codim1_indeterminacy}.  Their combination here is special to
group-variety targets: a rational map to a general complete variety can have
a nonempty codimension-two indeterminacy locus.

The extension theorem feeds both the Albanese construction discussed above
and the usual rigidity consequences for rational maps between abelian
varieties.  Over a non-algebraically-closed field the local extension
arguments remain available after suitable residue-field bookkeeping, but
descent of the resulting morphism is an additional problem.
